\QHeading{第1题　简便计算。}

\par\noindent\makebox[\linewidth][c]{
\begin{tikzpicture}[x=1bp,y=-1bp,color=black,draw=black,text=black]
\begin{scope}\clip (56.0,123.0) rectangle (448.0,277.33);
\hdhead{32}{运算律——信息窗2}
\hdsection{1}{145}{简便计算。}
\q{124}{193}{260}{13}{$25\times 18\times 4$}
\q{510}{193}{260}{13}{$125\times 53\times 8$}
\q{106}{241}{180}{13}{$=25\times $}
\q{238}{241}{180}{13}{$\times 18$}
\q{106}{299}{180}{13}{$=$}
\q{192}{299}{180}{13}{$\times 18$}
\q{106}{356}{180}{13}{$=$}
\q{488}{240}{180}{13}{$=125\times $}
\q{631}{240}{180}{13}{$\times 53$}
\q{488}{299}{180}{13}{$=$}
\q{573}{299}{180}{13}{$\times 53$}
\q{488}{355}{180}{13}{$=$}
\begin{scope}[scale=.6666667,draw=black,text=black]\draw[line width=.5pt] (174,233) rectangle (232,268);\end{scope}
\begin{scope}[scale=.6666667,draw=black,text=black]\draw[line width=.5pt] (126,291) rectangle (185,326);\end{scope}
\begin{scope}[scale=.6666667,draw=black,text=black]\draw[line width=.5pt] (126,349) rectangle (185,384);\end{scope}
\begin{scope}[scale=.6666667,draw=black,text=black]\draw[line width=.5pt] (569,232) rectangle (627,267);\end{scope}
\begin{scope}[scale=.6666667,draw=black,text=black]\draw[line width=.5pt] (509,291) rectangle (566,326);\end{scope}
\begin{scope}[scale=.6666667,draw=black,text=black]\draw[line width=.5pt] (509,348) rectangle (566,383);\end{scope}
\hdsection{2}{419}{用简便方法计算下面各题。}
\q{103}{467}{260}{13}{$25\times 4\times 17$}
\q{357}{467}{260}{13}{$125\times 36\times 8$}
\q{613}{467}{260}{13}{$25\times 8\times 4\times 25$}
\q{103}{636}{260}{13}{$360\div 45$}
\q{357}{636}{260}{13}{$25\times 16$}
\q{613}{636}{260}{13}{$125\times 32$}
\hdsection{3}{815}{解决问题。}
\q{104}{860}{341}{12}{① 会议厅里摆放绿植，一共摆放}
\q{104}{902}{341}{12}{　4层，每层有25盆绿植。会议}
\q{104}{944}{341}{12}{　厅里有2个这样的摆放区，一}
\q{104}{986}{341}{12}{　共能摆放多少盆绿植?}
\q{486}{860}{350}{12}{② 工地运水泥，每次运50袋，每}
\q{486}{902}{350}{12}{　天运2次，一共运18天。一}
\q{486}{944}{350}{12}{　共运水泥多少袋?}

\begin{scope}[color=answerred,draw=answerred,text=answerred]\ans{127.33}{160.00}{50.00}{10}{$\vphantom{0}4$}
\ans{93.33}{198.67}{50.00}{10}{$\vphantom{0}100$}
\ans{88.67}{237.33}{50.00}{10}{$\vphantom{0}1800$}
\ans{392.67}{160.00}{50.00}{10}{$\vphantom{0}8$}
\ans{343.33}{198.00}{50.00}{10}{$\vphantom{0}1000$}
\ans{342.00}{236.00}{50.00}{10}{$\vphantom{0}53000$}
\ans{68.67}{337.33}{166.67}{9.3}{$\setlength{\jot}{-2pt}\vphantom{0}\begin{aligned}[t]&=(25\times4)\times17\\&=100\times17\\&=1700\end{aligned}$}
\ans{238.67}{337.33}{166.67}{9.3}{$\setlength{\jot}{-2pt}\vphantom{0}\begin{aligned}[t]&=(125\times8)\times36\\&=1000\times36\\&=36000\end{aligned}$}
\ans{408.67}{337.33}{166.67}{9.3}{$\setlength{\jot}{-2pt}\vphantom{0}\begin{aligned}[t]&=(25\times4)\times(8\times25)\\&=100\times200\\&=20000\end{aligned}$}
\ans{68.67}{450.00}{166.67}{9.3}{$\setlength{\jot}{-2pt}\vphantom{0}\begin{aligned}[t]&=360\div9\div5\\&=40\div5\\&=8\end{aligned}$}
\ans{239.33}{450.00}{166.67}{9.3}{$\setlength{\jot}{-2pt}\vphantom{0}\begin{aligned}[t]&=25\times4\times4\\&=100\times4\\&=400\end{aligned}$}
\ans{408.67}{450.00}{166.67}{9.3}{$\setlength{\jot}{-2pt}\vphantom{0}\begin{aligned}[t]&=125\times8\times4\\&=1000\times4\\&=4000\end{aligned}$}
\ans{72.00}{692.00}{226.67}{10}{$\setlength{\jot}{-2pt}\begin{aligned}[t]&\phantom{={}}25\times4\times2\\&=100\times2\\&=200\text{（盆）}\end{aligned}$}
\ans{72.00}{740.00}{226.67}{10}{答：共能摆放200盆绿植。}
\ans{328.67}{664.00}{226.67}{10}{$\setlength{\jot}{-2pt}\begin{aligned}[t]&\phantom{={}}50\times2\times18\\&=100\times18\\&=1800\text{（袋）}\end{aligned}$}
\ans{328.67}{712.00}{226.67}{10}{答：共运水泥1800袋。}
\end{scope}
\end{scope}\pgfresetboundingbox\path[use as bounding box] (56.0,123.0) rectangle (448.0,277.33);
\end{tikzpicture}}\par
\par\noindent




\par



\par\noindent

\QHeading{第2题　用简便方法计算下面各题。}

\CalcBlock{1}{\begin{aligned}[t]&\phantom{={}}\color{black}25\times 4\times 17\\&\color{answerred}=(25\times4)\times17\\&\color{answerred}=100\times17\\&\color{answerred}=1700\end{aligned}}
\par

\CalcBlock[方法一：简便计算]{2}{\begin{aligned}[t]&\phantom{={}}\color{black}125\times 36\times 8\\&\color{answerred}=(125\times8)\times36\\&\color{answerred}=1000\times36\\&\color{answerred}=36000\end{aligned}}
\MethodBlock{方法二：结合因数，先凑整}{用乘法交换律、结合律，先算容易凑整的两个因数。}{\begin{aligned}[t]&\phantom{={}}125\times 36\times 8\\
&=(125\times 36)\times 8\\
&=4500\times 8\\
&=36000\end{aligned}}

\CalcBlock[方法一：简便计算]{3}{\begin{aligned}[t]&\phantom{={}}\color{black}25\times 8\times 4\times 25\\&\color{answerred}=(25\times4)\times(8\times25)\\&\color{answerred}=100\times200\\&\color{answerred}=20000\end{aligned}}
\MethodBlock{方法二：结合因数，先凑整}{用乘法交换律、结合律，先算容易凑整的两个因数。}{\begin{aligned}[t]&\phantom{={}}25\times 8\times 4\times 25\\
&=(4\times 25)\times 25\times 8\\
&=100\times 25\times 8\\
&=2500\times 8\\
&=20000\end{aligned}}

\MethodBlock{方法三：先处理括号或乘除，再拆因数，凑整计算}{把一个因数拆成两个因数，先凑成整百或整千。}{\begin{aligned}[t]&\phantom{={}}25\times 8\times 4\times 25\\
&=200\times 4\times 25\\
&=800\times 25\\
&=25\times (4\times 200)\\
&=(25\times 4)\times 200\\
&=100\times 200\\
&=20000\end{aligned}}

\MethodBlock{方法四：先处理括号或乘除，再把特殊因数改成整百或整千}{先把25、125等改写成整百或整千除以一个数。}{\begin{aligned}[t]&\phantom{={}}25\times 8\times 4\times 25\\
&=200\times 4\times 25\\
&=800\times 25\\
&=(100\div 4)\times 800\\
&=100\times (800\div 4)\\
&=100\times 200\\
&=20000\end{aligned}}

\CalcBlock[方法一：简便计算]{4}{\begin{aligned}[t]&\phantom{={}}\color{black}360\div 45\\&\color{answerred}=360\div (9\times 5)\\&\color{answerred}=360\div 9\div 5\\&\color{answerred}=40\div 5\\&\color{answerred}=8\end{aligned}}
\MethodBlock{方法二：拆分除数，连续相除}{}{\begin{aligned}[t]&\phantom{={}}360\div 45\\
&=360\div (5\times 9)\\
&=360\div 5\div 9\\
&=72\div 9\\
&=8\end{aligned}}

\MethodBlock{方法三：利用商不变的性质}{被除数和除数同时乘同一个非零数，商不变。}{\begin{aligned}[t]&\phantom{={}}360\div 45\\
&=(360\times 2)\div (45\times 2)\\
&=720\div 90\\
&=8\end{aligned}}

\CalcBlock[方法一：简便计算]{5}{\begin{aligned}[t]&\phantom{={}}\color{black}25\times 16\\&\color{answerred}=25\times (4\times 4)\\&\color{answerred}=(25\times 4)\times 4\\&\color{answerred}=100\times 4\\&\color{answerred}=400\end{aligned}}
\MethodBlock{方法二：把特殊因数改成整百或整千}{先把25、125等改写成整百或整千除以一个数。}{\begin{aligned}[t]&\phantom{={}}25\times 16\\
&=(100\div 4)\times 16\\
&=100\times (16\div 4)\\
&=100\times 4\\
&=400\end{aligned}}

\MethodBlock{方法三：按数位拆分计算}{把一个因数拆成整十、整百数与余下部分之和。}{\begin{aligned}[t]&\phantom{={}}25\times 16\\
&=25\times (10+6)\\
&=25\times 10+25\times 6\\
&=250+150\\
&=400\end{aligned}}

\CalcBlock[方法一：简便计算]{6}{\begin{aligned}[t]&\phantom{={}}\color{black}125\times 32\\&\color{answerred}=125\times (8\times 4)\\&\color{answerred}=(125\times 8)\times 4\\&\color{answerred}=1000\times 4\\&\color{answerred}=4000\end{aligned}}
\MethodBlock{方法二：把特殊因数改成整百或整千}{先把25、125等改写成整百或整千除以一个数。}{\begin{aligned}[t]&\phantom{={}}125\times 32\\
&=(1000\div 8)\times 32\\
&=1000\times (32\div 8)\\
&=1000\times 4\\
&=4000\end{aligned}}

\MethodBlock{方法三：凑整后补偿}{把接近整十、整百的因数凑整，再补回或减去多算的部分。}{\begin{aligned}[t]&\phantom{={}}125\times 32\\
&=125\times (30+2)\\
&=125\times 30+125\times 2\\
&=3750+250\\
&=4000\end{aligned}}

\MethodBlock{方法四：按数位拆分计算}{把一个因数拆成整十、整百数与余下部分之和。}{\begin{aligned}[t]&\phantom{={}}125\times 32\\
&=32\times (100+25)\\
&=32\times 100+32\times 25\\
&=3200+800\\
&=4000\end{aligned}}

\par\noindent

\QHeading{第3题　解决问题。}

\Problem{① 会议厅里摆放绿植，一共摆放4层，每层有25盆绿植。会议厅里有2个这样的摆放区，一共能摆放多少盆绿植?}
\par


\MethodTitle{方法一}
\Answer{\Formula{\begin{aligned}[t]&\phantom{={}}25\times4\times2\\&=100\times2\\&=200\text{（盆）}\end{aligned}}}
\Answer{答：共能摆放200盆绿植。}

\MethodTitle{方法二：先求两个摆放区共有多少层}
\Answer{每个摆放区4层，两个区8层，每层25盆。}
\Answer{\Formula{\begin{aligned}[t]&\phantom{={}}25\times (4\times 2)\\
&=25\times 8\\
&=200\text{（盆）}\end{aligned}}}
\Answer{答：一共能摆放200盆绿植。}
\Problem{② 工地运水泥，每次运50袋，每天运2次，一共运18天。一共运水泥多少袋?}

\MethodTitle{方法一}
\Answer{\Formula{\begin{aligned}[t]&\phantom{={}}50\times2\times18\\&=100\times18\\&=1800\text{（袋）}\end{aligned}}}
\Answer{答：共运水泥1800袋。}

\MethodTitle{方法二：先求一共运输多少次}
\Answer{每天运2次，18天共36次，每次50袋。}
\Answer{\Formula{\begin{aligned}[t]&\phantom{={}}50\times (2\times 18)\\
&=50\times 36\\
&=1800\text{（袋）}\end{aligned}}}
\Answer{答：一共运水泥1800袋。}
